\section{Giải phương trình \enquote{đạo hàm bằng 1}}
\label{sec:ch04-dong-loi-khong-doi}

\index{vòng quay lỗi không đổi!phương trình gốc}%
Mục 3.2 của bài bắt đầu từ chỗ nhỏ nhất có thể: một unit duy nhất, tự nối vào
chính nó bằng đúng một trọng số \(w_{jj}\), không có gì khác. Viết
\(\vartheta_j(t)\) cho tín hiệu lỗi tới unit \(j\) ở bước \(t\), tức ký hiệu
của bài cho đại lượng mà chương~\ref{ch:ky-hieu-va-bptt} đẩy ngược qua thời
gian. Với một unit như trên, tín hiệu ấy đi ngược một bước thì bị nhân với

\begin{equation}
  \pd{\vartheta_j(t)}{\vartheta_j(t+1)} = f_j'(a_j(t))\, w_{jj},
  \label{eq:ch04-thua-so-mot-unit}
\end{equation}

trong đó \(a_j(t)\) là \tn{tiền kích hoạt}{pre-activation} của unit và \(f_j\)
là \tn{hàm kích hoạt}{activation function} của nó. Bài không hỏi làm sao cho
thừa số này gần 1. Bài đặt nó bằng 1 rồi giải:

\begin{equation}
  f_j'(a_j(t))\, w_{jj} = 1.0 \quad \text{với mọi } t.
  \label{eq:ch04-rang-buoc}
\end{equation}

Đây là một phương trình vi phân theo \(a_j\), và nó chỉ có một họ nghiệm. Lấy
tích phân hai vế theo \(a_j(t)\):

\begin{equation}
  f_j(a_j(t)) = \frac{a_j(t)}{w_{jj}} + C.
  \label{eq:ch04-nghiem}
\end{equation}

Đọc kết quả này chậm, vì nó nói hai chuyện chứ không phải một. Chuyện thứ nhất:
\(f_j\) phải \emph{affine}, và nếu muốn unit giữ nguyên activation chứ không
trôi đều thì hằng số \(C\) phải bằng không, tức \(f_j\) tuyến tính. Không phải
\enquote{nên chọn hàm ít bão hòa}, mà là tuyến tính, không còn lựa chọn nào
khác. Bài in nghiệm với \(C = 0\) mà không nhắc tới hằng số; bài tập~1 cuối
chương hỏi lại chỗ này. Chuyện thứ hai đến khi thay \(w_{jj} = 1.0\) vào, lúc
đó \(f_j\) là ánh xạ đồng nhất và

\begin{equation}
  y_j(t+1) = f_j(w_{jj}\, y_j(t)) = y_j(t).
  \label{eq:ch04-giu-nguyen}
\end{equation}

Activation của unit đứng yên. Không phải nó nhớ lâu; nó không quên gì cả, vì
không có gì trong phép cập nhật làm nó thay đổi. Bài gọi cấu trúc này là
\tn{vòng quay lỗi không đổi}{constant error carousel}, viết tắt là CEC, và nó
là hạt nhân của toàn bộ kiến trúc còn lại.

\subsection*{Lệch khỏi 1 đắt bao nhiêu}

\index{vòng quay lỗi không đổi!vì sao cố định trọng số}%
Một câu hỏi hợp lý ở đây là vì sao phải \emph{cố định} \(w_{jj}\) ở 1 thay vì
để mạng học lấy một giá trị gần 1. Chương~\ref{ch:phan-tich} đã trả lời bằng lý
thuyết, nhưng con số cụ thể thì đáng nhìn. Với unit tuyến tính, dòng lỗi qua
\(q\) bước đúng bằng \(w^{q}\):

\begin{minted}{text}
     w  q=1         q=10        q=50        q=100
  0.90  9.0000e-01  3.4868e-01  5.1538e-03  2.6561e-05
  0.99  9.9000e-01  9.0438e-01  6.0501e-01  3.6603e-01
  1.00  1.0000e+00  1.0000e+00  1.0000e+00  1.0000e+00
  1.01  1.0100e+00  1.1046e+00  1.6446e+00  2.7048e+00
  1.10  1.1000e+00  2.5937e+00  1.1739e+02  1.3781e+04
\end{minted}

Lệch một phần trăm, tức \(w = 0.99\), làm tín hiệu còn 0.36603 sau 100 bước;
lệch một phần trăm theo chiều kia làm nó thành 2.7048. Lệch mười phần trăm cho
0.000026561 và 13781. Một tham số học được thì trong lúc huấn luyện nó \emph{sẽ}
lang thang quanh 1, và bảng trên là bảng giá của việc lang thang. Cố định nó
không phải là lười; đó là chỗ duy nhất trong kiến trúc mà một con số được biết
chính xác, và tiêu nó đi để đổi lấy một tham số nữa là một đổi chác tồi.

\subsection*{Vì sao một unit sigmoid không làm được}

\index{bão hòa!chặn 1/4 của sigmoid}%
Ràng buộc~\eqref{eq:ch04-rang-buoc} cũng loại luôn khả năng dùng một unit phi
tuyến thông thường, và đáng bỏ ra hai phút để thấy nó loại theo kiểu nào.

Với sigmoid, \(f'\) không bao giờ vượt \(1/4\). Vậy để tích \(f'(a) w\) chạm
tới 1 thì cần \(w \ge 4\), đó là \tn{điều kiện cần}{necessary condition} và nó
đến từ số học chứ không từ
huấn luyện. Nhưng đúng tại \(w = 4\) thì \(f'(a) \cdot 4 = 1\) chỉ đúng khi
\(f'(a) = 1/4\), tức chỉ tại đúng một điểm \(a = 0\). Trạng thái rời khỏi điểm
đó là ràng buộc gãy. Tôi cho chạy thử một unit sigmoid tự nối, không input,
không bias, 100 bước:

\begin{minted}{text}
     w  factor t=1   factor t=100 product q=100
  1.00  0.235004     0.224704     1.5382e-65
  2.00  0.393224     0.263401     2.1187e-58
  4.00  0.419974     0.076133     1.2052e-111
  8.00  0.141302     0.002689     5.5338e-256
\end{minted}

\begin{measured}{unit sigmoid không giữ được dòng lỗi ở bất kỳ trọng số nào}
\(w = 4\) là trọng số nhỏ nhất mà về lý thuyết ràng buộc có thể đạt được, và nó
cho tích \(1.2052\) nhân mười mũ trừ 111, tức \emph{tệ hơn} \(w = 2\) khoảng
năm mươi ba bậc độ lớn. Ở \(w = 8\) thì hệ số mỗi bước tụt xuống 0.002689 ngay
tại bước 100.

Đo bằng \code{experiments/ch04\_cec.py} tại tag \code{ch04}.
\end{measured}

Thứ tự của các con số ấy mới là chỗ đáng nhìn, chứ không phải độ nhỏ của chúng.
Trọng số lớn nhất trong bảng cho tích nhỏ nhất, và lý do là
\tn{bão hòa}{saturation}: \(w\) lớn đẩy tiền kích hoạt ra xa gốc, \(f'\) sụt
theo, và mức sụt ấy ăn hết phần mà \(w\) lớn mang lại. Không có \(w\) nào cứu
được unit sigmoid, và bảng trên cho thấy lý do không phải là \(w\) chưa đủ lớn.
Đường thoát duy nhất là đổi \(f\), và
phương trình~\eqref{eq:ch04-nghiem} đã nói đổi thành gì.

Đến đây bài đã có một unit giữ được thông tin vô hạn lâu. Nó cũng có một unit
vô dụng: một hằng số không đọc được input và không ai đọc được nó. Mục sau là
chỗ bài tự phá cái nghiệm nó vừa tìm ra.
